\hypertarget{appendix-d1ea57347893}{}
\section*{Dual quaternions and rigid motion}
\chapter{Dual quaternions and rigid motion}\label{appendix-d1ea57347893-dual-quaternions-and-rigid-motion}

\textbf{Curriculum depth:} Introduced · \textbf{First seen:} Phase7 · \textbf{Latest development:} Phase7

\section{Definition and mechanism}\label{appendix-d1ea57347893-definition-and-mechanism}

A rigid dual quaternion consists of a unit real quaternion and a dual quaternion part satisfying real dot dual equals zero. Encoding dual=(0,t)real/2 combines rotation and translation; composition follows the Hamilton convention. It does not represent arbitrary scale or shear.

\textbf{Reconstruction cue:} Recover translation with twice dual times conjugate(real); normalization must preserve the coupled constraint.

\section{Prerequisites}\label{appendix-d1ea57347893-prerequisites}

\begin{itemize}
\tightlist
\item
  \hyperref[appendix-965d9260eeab]{Quaternion algebra and Hamilton product}
\item
  \hyperref[appendix-0bf0ca73f735]{Linear affine and projective transformations}
\end{itemize}

\section{Used by}\label{appendix-d1ea57347893-used-by}

\begin{itemize}
\tightlist
\item
  \hyperref[appendix-292c061cd47c]{Transform hierarchies TRS and shear}
\end{itemize}

\section{Recall and next depth}\label{appendix-d1ea57347893-recall-and-next-depth}

Reconstruct the example, state its conventions and find an input that violates its preconditions. Deeper algorithms and measured production performance remain separate from the present model.

\section{Source and navigation}\label{appendix-d1ea57347893-source-and-navigation}

\href{Continuous\%20Notes/Computing\%20Foundations\%20-\%20From\%20Bits\%20to\%20C++\%20and\%20Systems\%20APIs\#Keep\%20the\%20promises\%20visible\%20while\%20moving\%20between\%20representations}{Continuous explanation} · \hyperref[appendix-fd090d982a15-phase-7---parts-601-700]{Coverage and source gap} · \hyperref[appendix-a103e7846d21-phase-7---source-analysis-and-verification]{Sources and verification}
